\subsection{Компьютерное зрение в контуре управления}

До сих пор изображение с камеры только показывалось. В этом разделе оно \emph{влияет на полёт}. Все три примера устроены по одной схеме замкнутого контура: камера $\to$ распознавание $\to$ ошибка $\to$ регулятор $\to$ команда скорости $\to$ движение дрона $\to$ новый кадр.

\begin{center}
\begin{tikzpicture}[node distance=5mm,
    b/.style={draw=GeoskanBlue, fill=GeoskanBlue!8, rounded corners=2pt, font=\sffamily\small, minimum height=8mm, align=center},
    a/.style={-latex, thick}]
\node[b] (cam) {Камера\\\footnotesize get\_cv\_frame};
\node[b, right=of cam] (cv) {Распознавание\\\footnotesize OpenCV / ArUco};
\node[b, right=of cv] (err) {Ошибка\\\footnotesize $e=$ цель $-$ факт};
\node[b, right=of err] (reg) {П-регулятор\\\footnotesize $u=K_p e$};
\node[b, below=8mm of cv, minimum width=60mm] (drone) {Дрон\\\footnotesize set\_manual\_speed\_body\_fixed};
\draw[a] (cam) -- (cv); \draw[a] (cv) -- (err); \draw[a] (err) -- (reg);
\draw[a] (reg) |- (drone); \draw[a] (drone) -| (cam);
\end{tikzpicture}
\end{center}

\begin{rulebox}[Безопасность контура]
Если кадра или цели нет, надо отправлять нулевую скорость: дрон зависнет, а не продолжит выполнять последнюю команду. Команды скорости короткие (\verb|interval=0.3|), коэффициенты начинайте с малых значений: слишком большой коэффициент $K_p$ раскачивает дрон --- он «проскакивает» цель и колеблется; уменьшите $K_p$ вдвое.
\end{rulebox}

\subsubsection*{Поворот к цветному объекту}

\begin{codefile}[python]{python-onboard/pioneer_sdk2/18_color_tracking_yaw.py}
\end{codefile}

\paragraph{Что происходит.} Цвет ищется в пространстве HSV: оттенок (H), насыщенность (S) и яркость (V) описывают цвет так, что порог меньше зависит от освещения. У красного два диапазона оттенка на краях шкалы (0--10 и 170--180), их маски объединяются операцией \verb/|=/. Морфологическое открытие \verb|MORPH_OPEN| удаляет одиночные шумовые пиксели. Центр масс пятна считается через моменты: $c_x = m_{10}/m_{00}$, $c_y = m_{01}/m_{00}$. П-регулятор: $\omega = K_p\,(w/2 - c_x)$, ограниченный \verb|MAX_YAW_RATE|; мёртвая зона \verb|DEAD_ZONE| убирает дрожание около центра. Знак: цель левее центра --- ошибка положительна --- \verb|yaw_rate > 0|, то есть поворот влево. Если кадра нет, отправляется нулевая скорость.

\begin{taskbox}[Другой цвет, высота и дистанция]
Подберите диапазон HSV для зелёного или синего предмета (подсказка: напишите программу-«пипетку», печатающую HSV центра кадра). Добавьте управление по высоте: держите объект по центру и по вертикали с помощью \verb|vz|. Добавьте дистанцию по площади пятна: \verb|vy| пропорционально $(S_\text{цель} - S)$.
\expected{дрон висит и плавно поворачивается вслед за предметом выбранного цвета; при потере предмета останавливается; поток \texttt{.../color/} показывает найденный центр.}
\end{taskbox}

\subsubsection*{ArUco-маркеры: распознавание и расстояние}

\begin{codefile}[python]{python-onboard/pioneer_sdk2/19_aruco_distance.py}
\end{codefile}

\paragraph{Что происходит.} Для расчёта нужна калибровка камеры: матрица камеры \verb|mtx| и коэффициенты искажений \verb|dist| хранятся в файле \texttt{data.yml}; его создаёт программа калибровки по шахматной доске \texttt{camera\_calibration.py} из репозитория примеров Pioneer-SDK2 (15 снимков доски $6\times9$ внутренних углов с разных ракурсов). \verb|ArucoDetector.detectMarkers()| возвращает углы, номера и отвергнутые кандидаты. \verb|solvePnP| решает обратную задачу: зная реальные 3D-координаты углов маркера (\verb|MARKER_POINTS|) и их положение на снимке, находит поворот \verb|rvec| и смещение \verb|tvec| маркера относительно камеры; флаг \verb|SOLVEPNP_IPPE_SQUARE| предназначен для плоских квадратов. Расстояние --- длина вектора $d = \lVert\mathbf t\rVert = \sqrt{t_x^2+t_y^2+t_z^2}$. \verb|drawFrameAxes| рисует оси маркера --- визуальная проверка, правильно ли решена задача.

\begin{warnbox}[Размер маркера]
Ошибка в \verb|MARKER_SIZE| пропорционально искажает все расстояния: если указать 0.10~м для маркера со стороной 0.05~м, дистанции окажутся вдвое больше. Измеряйте чёрный квадрат маркера линейкой.
\end{warnbox}

\subsubsection*{Удержание дистанции до маркера}

\begin{codefile}[python]{python-onboard/pioneer_sdk2/20_aruco_follow.py}
\end{codefile}

\paragraph{Что происходит.} Два потока: \verb|MarkerTracker| непрерывно читает камеру и решает PnP, основной поток раз в 0.1~с отправляет команды скорости --- медленная обработка кадра не задерживает управление. \verb|threading.Lock| защищает общие данные (\verb|tvec|, \verb|seen_at|) от одновременного чтения и записи. Потеря цели: \verb|latest()| возвращает \verb|None|, если маркер не виден дольше \verb|LOST_TIMEOUT|, и тогда скорость нулевая. В осях камеры $z$ --- расстояние вперёд, $x$ --- вправо, регулятор:
\[ v_y = K_d\,(z - d^{*}),\qquad \omega = -K_\psi \operatorname{atan2}(x, z), \]
обе величины ограничены через \verb|np.clip|. Минус перед вторым коэффициентом: положительный \verb|yaw_rate| поворачивает влево.

\begin{taskbox}[Улучшить слежение]
Добавьте управление высотой по $t_y$, чтобы маркер оставался в центре кадра по вертикали. Добавьте интегральную составляющую (ПИ-регулятор) и сравните точность удержания дистанции. Сделайте посадку на маркер: снижайтесь, пока маркер виден под камерой (камера $-80^\circ$), затем \verb|land()|.
\expected{дрон держится около 1.2~м от маркера; при отходе маркера следует за ним; при потере маркера останавливается; посадка на маркер заканчивается \texttt{land()} над ним.}
\end{taskbox}

\begin{infobox}[Что запускается в симуляторе]
Примеры 18--20 используют камеру и OpenCV. Пример 18 можно пробовать в симуляторе, если на сцене есть яркий красный предмет и камера Мини~2 отдаёт кадры. Примеры 19 и 20 требуют файл калибровки \texttt{data.yml} и модуль \texttt{cv2.aruco}: их надёжно запускать только на дроне.
\end{infobox}
